%%-----------Chapter 3------------------------------------------
\chapter{Motion}

\section{What is Motion?}

\begin{keybox}[Motion and Rest]
The apparent change in a body's (or object's) position with time is called \textbf{motion}.
If the body does not change its position as time passes, it is said to be at \textbf{rest}.
\end{keybox}

\subsection{The relative nature of motion}

But when do we say that a body is not changing its position? A cup on the table remains on
the table --- we say the cup is at rest. However, if we station ourselves on Mars, the whole
Earth is found to be changing its position, and so the room, the table, and the cup are all
continuously changing their positions.
\begin{itemize}
  \item The cup is \emph{at rest} if viewed from the room.
  \item The cup is \emph{moving} if viewed from Mars.
\end{itemize}
Motion is therefore a combined property of the object under study and the observer. There
is no meaning of rest or motion without the viewer.

\begin{keybox}[Key Principle]
\centering
Nothing is in absolute rest or in absolute motion.\\
Mars is moving with respect to the cup; the cup moves with respect to Mars.
\end{keybox}

\begin{workedexample}[3.1: Frames and relative motion]
A passenger sits still inside a moving train.

(a) From the passenger's own frame: the passenger is \emph{at rest}.

(b) From a bystander on the platform: the passenger is \emph{moving} at the speed of the train.

(c) From an observer on the Moon: both the passenger and the bystander are moving (the
Earth rotates and orbits the Sun).

All three descriptions are equally valid --- rest and motion depend entirely on the reference
frame of the observer.
\end{workedexample}

\subsection{Frame of reference}

To locate the position of a particle we need a \textbf{frame of reference} --- a choice of a
spatial origin and coordinate axes to label position. We choose Cartesian coordinates, so the
coordinates $(x,y,z)$ of the particle $P$ specify its position relative to that frame. A clock is
added to the frame to measure time.

\begin{keybox}[Rest and motion in terms of coordinates]
\textbf{Rest:} If all three coordinates $x,y,z$ of the particle remain unchanged as time passes,
the particle is at rest with respect to this frame.

\textbf{Motion:} If any one or more coordinates change with time, the body is moving with
respect to this frame.
\end{keybox}

\begin{keybox}[One-dimensional motion first]
For our first look at this subject we will forget about the three dimensions of the world
and concentrate on moving in one direction --- like a car on a single straight road. We
will return to three dimensions after we see how to describe motion in one dimension.
\end{keybox}


\section{Velocities and the Derivative}

\subsection{Position function}

Suppose an object moves along a straight line. We describe its motion by an \textbf{equation of motion}
\[
  s = f(t),
\]
where $s$ is the displacement (direct distance from the origin along the line) of the object
at time $t$. The function $f$ that describes the motion is called the \textbf{position function}
of the object.

\subsection{Average velocity}

In the time interval from $t$ to $t+\Delta t$, the change in position (displacement) is
$f(t+\Delta t)-f(t)$.

\begin{keybox}[Average velocity]
The average velocity over the time interval $[t,\,t+\Delta t]$ is
\[
  v_\text{av} = \frac{\text{displacement}}{\text{time}} = \frac{f(t+\Delta t)-f(t)}{\Delta t}.
\]
Geometrically, this is the slope of the secant line $PQ$ on the position-time graph.
\end{keybox}

\begin{workedexample}[3.2: Computing average velocity]
A car moves along a straight road with position function $f(t)=3t^2+2t$ metres (where
$t$ is in seconds).

(a) Displacement from $t=1\,\text{s}$ to $t=3\,\text{s}$:
\[
  f(1)=3(1)^2+2(1)=5\;\text{m}, \qquad f(3)=3(3)^2+2(3)=33\;\text{m}, \qquad \Delta s=28\;\text{m}.
\]

(b) Average velocity over $[1,3]$:
\[
  v_\text{av} = \frac{f(3)-f(1)}{3-1} = \frac{28}{2} = 14\;\text{m/s}.
\]
\end{workedexample}

\begin{tryit}[3.1: Average velocity from a table]
A cyclist's position at various times is recorded below. Find the average velocity over
the full 15\,s interval.
\begin{center}
\begin{tabular}{ccccc}
  \toprule
  $t$ (s) & 0 & 5 & 10 & 15 \\
  \midrule
  $s$ (m) & 0 & 12 & 20 & 24 \\
  \bottomrule
\end{tabular}
\end{center}
\textit{(Answer: $v_\text{av}=24/15=1.6\;\text{m/s}$.)}
\end{tryit}

\subsection{Instantaneous velocity}

Now suppose we compute the average velocities over shorter and shorter time intervals
$[t,t+\Delta t]$ --- in other words, we let $\Delta t$ approach zero.

\begin{keybox}[Instantaneous velocity]
The instantaneous velocity $v(t)$ at time $t$ is defined as the limit of these average velocities:
\[
  v(t) = \lim_{\Delta t\to 0}\frac{f(t+\Delta t)-f(t)}{\Delta t} = \frac{df}{dt},
\]
provided that this limit exists. Geometrically, $v(t)$ equals the slope of the tangent line
$\ell$ to the position-time curve at the point $P$.
\end{keybox}

\begin{workedexample}[3.3: Instantaneous velocity from the limit definition]
Find the instantaneous velocity $v(t)$ of the car $f(t)=3t^2+2t$ from the limit definition.
Hence find $v(1)$ and $v(3)$.

\medskip\textit{Solution.} Expand $f(t+\Delta t)$:
\[
  f(t+\Delta t)=3(t+\Delta t)^2+2(t+\Delta t)=3t^2+6t\,\Delta t+3(\Delta t)^2+2t+2\,\Delta t.
\]
Subtract $f(t)=3t^2+2t$:
\[
  f(t+\Delta t)-f(t) = 6t\,\Delta t+3(\Delta t)^2+2\,\Delta t = \Delta t\,(6t+3\,\Delta t+2).
\]
Divide by $\Delta t$ and take $\Delta t\to 0$:
\[
  v(t) = \lim_{\Delta t\to 0}(6t+3\,\Delta t+2) = 6t+2\;\text{m/s}.
\]
At $t=1\,\text{s}$: $v(1)=6(1)+2=8\;\text{m/s}$. At $t=3\,\text{s}$: $v(3)=6(3)+2=20\;\text{m/s}$.

Note: the average velocity over $[1,3]$ was 14\,m/s, which lies between 8 and 20\,m/s --- exactly as expected!
\end{workedexample}

\begin{tryit}[3.2: Limit definition for $f(t)=t^2$]
Use the limit definition to show that for $f(t)=t^2$, the instantaneous velocity is
$v(t)=2t\;\text{m/s}$. Hence find the velocity at $t=3\,\text{s}$.

\textit{Hint:} $(t+\Delta t)^2-t^2=2t\,\Delta t+(\Delta t)^2$. \textit{Answer:} $v(3)=6\;\text{m/s}$.
\end{tryit}

\subsection{The derivative}

The quantity $df/dt$ found above is called the \textbf{derivative} of $f$ with respect to $t$.
The process of finding it is called \textbf{differentiating}.

\begin{keybox}[Tangent line]
The tangent line to $y=f(t)$ at the point $(t,f(t))$ is the unique line through $(t,f(t))$
whose slope equals $\dfrac{df(t)}{dt}$, the derivative of $f$ at $t$.
\end{keybox}

\subsection{Other notations for the derivative}

Let $y=f(x)$. All of the following mean the same thing:
\[
  \frac{dy}{dx} = \frac{df}{dx} = \frac{d}{dx}f(x) = f'(x) = y'.
\]
For derivatives with respect to time $t$, Newton used an overdot:
\[
  \dot{f} = \frac{df}{dt}, \qquad \dot{y} = \frac{dy}{dt}.
\]

\begin{center}
\begin{tabular}{ll}
  \toprule
  \textbf{Notation} & \textbf{Name / origin} \\
  \midrule
  $\dfrac{dy}{dx},\ \dfrac{df}{dx},\ \dfrac{d}{dx}f(x)$ & Leibniz notation \\[8pt]
  $f'(x),\ y'$ & Lagrange (prime) notation \\[4pt]
  $\dot{y},\ \dot{f}$ & Newton (dot) notation --- time derivatives \\
  \bottomrule
\end{tabular}
\end{center}

All notations are in common use. In physics, $\dot{s}=v$ for velocity is popular for time
derivatives. You should be comfortable reading all of them.


\section*{Formula Summary: Motion}
\addcontentsline{toc}{section}{Formula Summary: Motion}

\begin{formulabox}[Formula Summary: Motion]
\begin{tabular}{ll}
  \toprule
  \textbf{Quantity} & \textbf{Formula} \\
  \midrule
  Position function & $s=f(t)$ \\[4pt]
  Displacement over $[t,\,t+\Delta t]$ & $\Delta s = f(t+\Delta t)-f(t)$ \\[4pt]
  Average velocity & $v_\text{av} = \dfrac{f(t+\Delta t)-f(t)}{\Delta t}$ \\[8pt]
  Instantaneous velocity & $v(t) = \displaystyle\lim_{\Delta t\to 0}\frac{f(t+\Delta t)-f(t)}{\Delta t} = \dfrac{df}{dt}$ \\[8pt]
  Slope of tangent & $\dfrac{df}{dt}$ at that point \\
  \bottomrule
\end{tabular}
\end{formulabox}

\medskip
\textbf{An important note on vectors:} Because we are currently studying motion along a
single straight line (1D), the direction of our vectors is indicated simply by a positive or
negative sign. To keep the focus on the calculus, we will temporarily drop the vector notation
($\vec{v}$) and use scalars ($v$) to represent these 1D quantities.


\section{Velocity as a Derivative and Graphing}

\textbf{Ex 1. A falling ball ($s=f(t)=5t^2$).}

\begin{center}
\begin{tabular}{cc}
  \toprule
  $t$ (s) & $s$ (m) \\
  \midrule
  0 & 0 \\
  1 & 5 \\
  2 & 20 \\
  3 & 45 \\
  4 & 80 \\
  5 & 125 \\
  \bottomrule
\end{tabular}
\end{center}

The formula for this curve can be written as $s = f(t) = 5t^2$, which enables us to calculate
the distance at any time.

If we want to know how fast a ball is falling at exactly $t=4\,\text{s}$, we cannot simply
use a wide interval. Between 3 and 4 seconds, the ball averages 35\,m/s
($v_\text{av}=\frac{80-45}{4-3}$). However, it is constantly accelerating during that time.
That average does not reflect the true velocity at the exact 4-second mark.

To get closer to the truth, we look at a very small interval immediately following $t=4\,\text{s}$:
\begin{itemize}
  \item At 4.1\,s: total displacement $s=5(4.1)^2=84.05\,\text{m}$.
  \item At 4\,s: total displacement was 80\,m.
  \item Average velocity for this tiny window: $4.05/0.1=40.5\,\text{m/s}$.
\end{itemize}

Taking the general approach, if at time $t_0=4\,\text{s}$ the position is $s_0=5t_0^2$, then at
$t_0+\Delta t$ the new position is
\[
  s_0+\Delta s = 5(t_0+\Delta t)^2 = 5t_0^2+10t_0\,\Delta t+5\,\Delta t^2.
\]
So the extra displacement is $\Delta s = 10t_0\,\Delta t+5\,\Delta t^2$, and the average velocity is
\begin{equation}
  v_\text{av} = \frac{\Delta s}{\Delta t} = 10t_0+5\,\Delta t.
\end{equation}
Taking $\Delta t\to 0$:
\[
  v(\text{at time }t_0) = 10t_0.
\]
At $t_0=4\,\text{s}$: $v(4)=10\times4=40\,\text{m/s}$.

The procedure we have just carried out is performed so often in mathematics that special notation
has been assigned. The velocity is
\begin{equation}
  v = \lim_{\Delta t\to 0}\frac{\Delta s}{\Delta t} = \frac{ds}{dt}.
\end{equation}
The quantity $ds/dt$ is called the ``derivative of $s$ with respect to $t$''. For our example:
\[
  \frac{d}{dt}(5t^2) = 10t.
\]

\textbf{A Short Table of Derivatives.}

\begin{center}
\begin{tabular}{lc}
  \toprule
  \multicolumn{2}{l}{\textit{$u$, $v$, $w$ are arbitrary functions of $t$; $c$ and $n$ are arbitrary constants.}} \\
  \midrule
  \textbf{Function, $f(t)$} & \textbf{Derivative, $\dfrac{df}{dt}$} \\
  \midrule
  $c$ (constant) & $0$ \\[4pt]
  $t^n$ & $n\,t^{n-1}$ \\[4pt]
  $cu$ & $c\,\dfrac{du}{dt}$ \\[4pt]
  $\sin(t)$ & $\cos(t)$ \\[4pt]
  $\cos(t)$ & $-\sin(t)$ \\[4pt]
  $e^t$ & $e^t$ \\[4pt]
  $\ln(t)$ & $\dfrac{1}{t}$ \\[4pt]
  $u+v+w+\cdots$ & $\dfrac{du}{dt}+\dfrac{dv}{dt}+\dfrac{dw}{dt}+\cdots$ \\[4pt]
  $uv$ & $u\dfrac{dv}{dt}+v\dfrac{du}{dt}$ \\[4pt]
  $\dfrac{u}{v}$ & $\dfrac{v\,(du/dt)-u\,(dv/dt)}{v^2}$ \\
  \bottomrule
\end{tabular}
\end{center}

\textbf{Exercise.} Find the derivative of the following functions.
\begin{enumerate}[label=\alph*.]
  \item $f(t) = 3t - 8$
  \item $f(x) = mx + b$
  \item $f(t) = 2.5t^2 + 6t$
  \item $f(x) = 4 + 8x - 5x^2$
  \item $A(p) = 4p^3 + \ln(p)$
  \item $F(t) = e^t(t^2-1)$
  \item $f(x) = \dfrac{1}{x^2-4}$
  \item $F(v) = \dfrac{v}{v+2}$
  \item $g(u) = \dfrac{u+1}{4u-1}$
  \item $f(x) = x^4 + \sin(x)$
  \item $g(x) = \dfrac{1}{1+\sqrt{x}}$
\end{enumerate}

\textit{Example solutions:}
\[
  \text{(a)}\quad f'(t) = 3, \qquad
  \text{(d)}\quad f'(x) = 8 - 10x.
\]


\section{Distance as an Integral}

Now we discuss the inverse problem. Suppose that instead of a table of distances, we have a
table of speeds at different times, starting from zero. From our earlier result, $s=5t^2$ and
$v=10t$.

The distance travelled is equal to the area under the velocity--time curve. In general, when
the velocity--time graph is not a simple geometric shape, we break the motion into small
intervals $\Delta t$ and approximate the distance in each interval as $\Delta s\approx v\,\Delta t$.

In general, suppose an object moves with velocity $v=v(t)$, where $a\le t\le b$. We take
velocity readings at times $t_0(=a),t_1,t_2,\ldots,t_n(=b)$ so that the velocity is
approximately constant on each sub-interval with $\Delta t=(b-a)/n$. The total distance
is approximately
\[
  v(t_0)\Delta t + v(t_1)\Delta t + \cdots + v(t_{n-1})\Delta t = \sum_{i=0}^{n-1}v(t_i)\,\Delta t.
\]
The true distance is the limit as $\Delta t\to 0$ (equivalently, $n\to\infty$):
\begin{equation}
  s = \lim_{\Delta t\to 0}\sum_{i=0}^{n-1}v(t_i)\,\Delta t = \lim_{n\to\infty}\sum_{i=0}^{n-1}v(t_i)\,\Delta t.
\end{equation}
The mathematicians have invented a symbol for this limit. The $\Delta$ turns into a $d$, the
addition is written with a great ``S'' (from the Latin \emph{summa}), distorted into the integral sign:
\begin{equation}
  s = \int v(t)\,dt.
\end{equation}
This process of adding all these terms together is called \textbf{integration}, and it is the
opposite process to differentiation.

\textbf{A Short Table of Derivatives and Integrals.}

\begin{center}
\begin{tabular}{lcc}
  \toprule
  \multicolumn{3}{l}{\textit{$u$, $v$, $w$ are arbitrary functions of $t$; $c$, $k$, $n$ are arbitrary constants.}} \\
  \midrule
  \textbf{Function, $f(t)$} & \textbf{Derivative, $\dfrac{df}{dt}$} & \textbf{Integral, $\displaystyle\int f(t)\,dt$} \\
  \midrule
  $c$ & $0$ & $ct+k$ \\[6pt]
  $t^n$ & $n\,t^{n-1}$ & $\dfrac{t^{n+1}}{n+1}+k \quad (n\ne -1)$ \\[6pt]
  $cu$ & $c\,\dfrac{du}{dt}$ & $c\displaystyle\int u\,dt$ \\[6pt]
  $\sin(t)$ & $\cos(t)$ & $-\cos(t)+k$ \\[4pt]
  $\cos(t)$ & $-\sin(t)$ & $\sin(t)+k$ \\[4pt]
  $e^t$ & $e^t$ & $e^t+k$ \\[4pt]
  $u+v+w+\cdots$ & $\dfrac{du}{dt}+\dfrac{dv}{dt}+\cdots$ &
    $\displaystyle\int u\,dt+\int v\,dt+\cdots$ \\
  \bottomrule
\end{tabular}
\end{center}

\medskip
Returning to our falling ball example, $v(t)=10t$:
\[
  s = \int v(t)\,dt = \int 10t\,dt = 5t^2 + C.
\]
The distance travelled between $t=0$ and $t=5\,\text{s}$:
\[
  s = \int_0^5 v(t)\,dt = \int_0^5 10t\,dt = \bigl[5t^2\bigr]_0^5
    = 5(5)^2 - 5(0)^2 = 125\;\text{m},
\]
which agrees with the area under the velocity--time graph over $[0,5]$.

\textbf{Exercise.} Evaluate the integral.
\begin{enumerate}[label=\alph*.]
  \item $\displaystyle\int(3t-8)\,dt$
  \item $\displaystyle\int(mx+b)\,dx$
  \item $\displaystyle\int_0^3(5t^2+6t)\,dt$
  \item $\displaystyle\int_\pi^4(4+8x-5x^2)\,dx$
  \item $\displaystyle\int_1^2\bigl(3e^t+(t^2-1)\bigr)\,dt$
  \item $\displaystyle\int_0^\pi(x^4+\sin x)\,dx$
\end{enumerate}

\textit{Example solutions:}
\[
  \text{(a)}\;\int(3t-8)\,dt = \tfrac{3}{2}t^2 - 8t + k, \qquad
  \text{(c)}\;\int_0^3(5t^2+6t)\,dt = \frac{5}{3}(3^3-0^3)+3(3^2-0^2) = 45+27 = 72.
\]


\section{Acceleration}

If the velocity of a particle remains constant as time passes, we say that it is moving with
\textbf{uniform velocity}. If the velocity changes with time, the particle is said to be
\textbf{accelerating}.

The \textbf{average acceleration} over $[t_1,t_2]$ is
\begin{equation}
  a_\text{av} = \frac{v_2-v_1}{t_2-t_1} = \frac{\Delta v}{\Delta t}.
\end{equation}
The \textbf{instantaneous acceleration} is
\begin{equation}
  a(t) = \lim_{\Delta t\to 0}\frac{\Delta v}{\Delta t} = \frac{dv}{dt}.
\end{equation}
Since velocity is $ds/dt$, and acceleration is the time derivative of velocity:
\begin{equation}
  a(t) = \frac{dv}{dt} = \frac{d}{dt}\!\left(\frac{ds}{dt}\right) = \frac{d^2s}{dt^2},
\end{equation}
which is the \textbf{second derivative} of position with respect to time.

For the falling ball $s(t)=5t^2$:
\[
  v(t) = \frac{ds}{dt} = 10t, \qquad a(t) = \frac{dv}{dt} = 10\;\text{m/s}^2.
\]
The acceleration is constant because the gravitational force on the falling body is constant,
and by Newton's second law, acceleration is proportional to the net force.


\section{Motion with Constant Acceleration}

Suppose a particle moves along a straight line (taken as the $x$-axis) with constant
acceleration $a$. Let the velocity at $t=0$ be $v_0$ and at time $t$ be $v$. Then:
\[
  \frac{dv}{dt} = a \implies dv = a\,dt.
\]
Integrating from 0 to $t$ (velocity from $v_0$ to $v$):
\[
  \int_{v_0}^v dv = \int_0^t a\,dt \implies v - v_0 = at,
\]
\begin{equation}
  \boxed{v = v_0 + at.}
\end{equation}
Writing $dx/dt=v_0+at$ and integrating from $x_0$ to $x$ (time from 0 to $t$):
\[
  x - x_0 = v_0 t + \frac{1}{2}at^2,
\]
\begin{equation}
  \boxed{\Delta x = v_0 t + \tfrac{1}{2}at^2.}
\end{equation}
Squaring equation (10): $v^2=(v_0+at)^2=v_0^2+2av_0 t+a^2t^2=v_0^2+2a\bigl(v_0 t+\frac{1}{2}at^2\bigr)$, so
\begin{equation}
  \boxed{v^2 = v_0^2 + 2a\,\Delta x.}
\end{equation}

\begin{keybox}[Equations of Kinematics]
\[
  v = v_0 + at, \qquad
  \Delta x = v_0 t + \tfrac{1}{2}at^2, \qquad
  v^2 = v_0^2 + 2a\,\Delta x.
\]
\end{keybox}

The three equations above are called the \textbf{kinematic equations}. Remember that $x$
represents the position (or displacement) at time $t$ and not necessarily the total distance
travelled. The quantities $v_0$, $v$, and $a$ may take positive or negative values depending
on their direction.


\section{Freely Falling Bodies}

A common example of motion in a straight line with constant acceleration is the \textbf{free fall}
of a body near the earth's surface. If air resistance is neglected and a body is dropped, it
falls along a vertical straight line. The acceleration in the vertically downward direction has
magnitude approximately $9.8\,\text{m/s}^2$ and is known as the \textbf{acceleration due to
gravity}, denoted by $g$.

The sign of $g$ depends on the chosen positive direction. Taking vertically upward as
positive gives $a=-g$; taking downward as positive gives $a=+g$. Substituting into the
kinematic equations accordingly:
\[
  v = v_0 \pm g\,t, \qquad \Delta y = v_0 t \pm \tfrac{1}{2}g\,t^2, \qquad v^2 = v_0^2 \pm 2g\,\Delta y,
\]
where the upper sign ($+$) applies when downward is taken as positive, and the lower sign
($-$) when upward is taken as positive.

\textbf{Motion in a Plane (2D) and in Space (3D).}
In three dimensions the position vector is
\[
  \vec{s} = x\hat{x} + y\hat{y} + z\hat{z}.
\]
Differentiating with respect to time gives the velocity vector:
\[
  \vec{v} = \frac{d\vec{s}}{dt} = \frac{dx}{dt}\hat{x} + \frac{dy}{dt}\hat{y} + \frac{dz}{dt}\hat{z}
           = v_x\hat{x} + v_y\hat{y} + v_z\hat{z},
\]
with magnitude (speed)
\[
  v = |\vec{v}| = \sqrt{v_x^2+v_y^2+v_z^2}.
\]
Differentiating once more gives the acceleration vector:
\[
  \vec{a} = \frac{d\vec{v}}{dt} = a_x\hat{x} + a_y\hat{y} + a_z\hat{z}, \qquad
  a = \sqrt{a_x^2+a_y^2+a_z^2}.
\]


\section{Projectile Motion}

An important example of motion in a plane with constant acceleration is \textbf{projectile
motion}. When a particle is thrown obliquely near the earth's surface, it moves along a
curved path. Such a particle is called a \textbf{projectile}. We assume the particle remains
close to the surface and air resistance is negligible, so acceleration is constant in the
vertically downward direction with magnitude $g\approx9.8\,\text{m/s}^2$.

A particle is projected from point $O$ with initial speed $v_0$ at angle $\theta$ with the horizontal.
The initial conditions are:
\begin{equation}
  x_0=0, \quad v_{0,x}=v_0\cos\theta, \quad a_x=0,
\end{equation}
\begin{equation}
  y_0=0, \quad v_{0,y}=v_0\sin\theta, \quad a_y=-g.
\end{equation}

\textbf{Horizontal Motion} ($a_x=0$):
\begin{equation}
  v_x = v_0\cos\theta \qquad \text{(constant)},
\end{equation}
\begin{equation}
  x = v_0 t\cos\theta.
\end{equation}

\textbf{Vertical Motion:}
\begin{equation}
  v_y = v_0\sin\theta - g\,t,
\end{equation}
\begin{equation}
  y = v_0 t\sin\theta - \tfrac{1}{2}g\,t^2,
\end{equation}
\begin{equation}
  v_y^2 = (v_0\sin\theta)^2 - 2g\,y.
\end{equation}

\textbf{Time of Flight.} At point $B$ (landing), $y=0$:
\[
  0 = v_0 t\sin\theta - \tfrac{1}{2}g\,t^2 = t\bigl(v_0\sin\theta - \tfrac{1}{2}g\,t\bigr).
\]
The non-trivial solution gives the time of flight:
\begin{equation}
  T = \frac{2v_0\sin\theta}{g}.
\end{equation}

\textbf{Range.} The horizontal range $OB$ is the distance travelled during time $T$:
\begin{equation}
  OB = v_0\,T\cos\theta = \frac{2v_0^2\sin\theta\cos\theta}{g} = \frac{v_0^2\sin2\theta}{g}.
\end{equation}
Notice that the greatest range is achieved if the projectile is thrown at $45^\circ$ to the
horizontal.

\textbf{Maximum Height Reached.} At the peak, $v_y=0$. From (18): $t=v_0\sin\theta/g$. Substituting into (19):
\begin{equation}
  H = \frac{v_0^2\sin^2\theta}{2g}.
\end{equation}
The time to reach maximum height is exactly half the total time of flight.

\textbf{Equation of Trajectory.} Eliminating $t$ between (16) and (19):
\begin{equation}
  y = x\tan\theta - \frac{g\,x^2}{2v_0^2\cos^2\theta}.
\end{equation}
This is the equation of the trajectory --- a \textbf{parabola}.
